% Options for packages loaded elsewhere
\PassOptionsToPackage{unicode}{hyperref}
\PassOptionsToPackage{hyphens}{url}
%
\documentclass[
  12pt,
  a4paper,
]{article}
\usepackage{amsmath,amssymb}
\usepackage{iftex}
\ifPDFTeX
  \usepackage[T1]{fontenc}
  \usepackage[utf8]{inputenc}
  \usepackage{textcomp} % provide euro and other symbols
\else % if luatex or xetex
  \usepackage{unicode-math} % this also loads fontspec
  \defaultfontfeatures{Scale=MatchLowercase}
  \defaultfontfeatures[\rmfamily]{Ligatures=TeX,Scale=1}
\fi
\usepackage{lmodern}
\ifPDFTeX\else
  % xetex/luatex font selection
\fi
% Use upquote if available, for straight quotes in verbatim environments
\IfFileExists{upquote.sty}{\usepackage{upquote}}{}
\IfFileExists{microtype.sty}{% use microtype if available
  \usepackage[]{microtype}
  \UseMicrotypeSet[protrusion]{basicmath} % disable protrusion for tt fonts
}{}
\makeatletter
\@ifundefined{KOMAClassName}{% if non-KOMA class
  \IfFileExists{parskip.sty}{%
    \usepackage{parskip}
  }{% else
    \setlength{\parindent}{0pt}
    \setlength{\parskip}{6pt plus 2pt minus 1pt}}
}{% if KOMA class
  \KOMAoptions{parskip=half}}
\makeatother
\usepackage{xcolor}
\usepackage[margin=1in]{geometry}
\usepackage{longtable,booktabs,array}
\usepackage{calc} % for calculating minipage widths
% Correct order of tables after \paragraph or \subparagraph
\usepackage{etoolbox}
\makeatletter
\patchcmd\longtable{\par}{\if@noskipsec\mbox{}\fi\par}{}{}
\makeatother
% Allow footnotes in longtable head/foot
\IfFileExists{footnotehyper.sty}{\usepackage{footnotehyper}}{\usepackage{footnote}}
\makesavenoteenv{longtable}
\setlength{\emergencystretch}{3em} % prevent overfull lines
\providecommand{\tightlist}{%
  \setlength{\itemsep}{0pt}\setlength{\parskip}{0pt}}
\setcounter{secnumdepth}{-\maxdimen} % remove section numbering
\newlength{\cslhangindent}
\setlength{\cslhangindent}{1.5em}
\newlength{\csllabelwidth}
\setlength{\csllabelwidth}{3em}
\newlength{\cslentryspacingunit} % times entry-spacing
\setlength{\cslentryspacingunit}{\parskip}
\newenvironment{CSLReferences}[2] % #1 hanging-ident, #2 entry spacing
 {% don't indent paragraphs
  \setlength{\parindent}{0pt}
  % turn on hanging indent if param 1 is 1
  \ifodd #1
  \let\oldpar\par
  \def\par{\hangindent=\cslhangindent\oldpar}
  \fi
  % set entry spacing
  \setlength{\parskip}{#2\cslentryspacingunit}
 }%
 {}
\usepackage{calc}
\newcommand{\CSLBlock}[1]{#1\hfill\break}
\newcommand{\CSLLeftMargin}[1]{\parbox[t]{\csllabelwidth}{#1}}
\newcommand{\CSLRightInline}[1]{\parbox[t]{\linewidth - \csllabelwidth}{#1}\break}
\newcommand{\CSLIndent}[1]{\hspace{\cslhangindent}#1}
\ifdefined\LTcapwidth\setlength{\LTcapwidth}{\textwidth}\fi
\ifLuaTeX
  \usepackage{selnolig}  % disable illegal ligatures
\fi
\IfFileExists{bookmark.sty}{\usepackage{bookmark}}{\usepackage{hyperref}}
\IfFileExists{xurl.sty}{\usepackage{xurl}}{} % add URL line breaks if available
\urlstyle{same}
\hypersetup{
  hidelinks,
  pdfcreator={LaTeX via pandoc}}

\author{}
\date{}

\begin{document}

\hypertarget{digital-twins-in-manufacturing-a-survey}{%
\section{Digital Twins in Manufacturing: A
Survey}\label{digital-twins-in-manufacturing-a-survey}}

This survey maps the literature on digital twins as they apply to
manufacturing. It covers what a digital twin is and how it is
distinguished from related concepts; how the twin is kept synchronised
with its physical counterpart; how twins are deployed on the factory
floor; and the interoperability standards that determine whether a twin
project survives integration. It does not cover machine-learning anomaly
detection except where a source ties it to a twin. The reader is assumed
to be familiar with manufacturing systems at a general level and to need
a structured map of the twin literature rather than an introduction to
manufacturing itself.

\hypertarget{definitions-and-a-three-level-taxonomy}{%
\subsection{Definitions and a Three-Level
Taxonomy}\label{definitions-and-a-three-level-taxonomy}}

The term \emph{digital twin} is used loosely across the literature, and
much of the early literature disagrees about where a simulation ends and
a twin begins. Author and Builder {[}1{]} resolve this by fixing the
vocabulary around integration depth. They distinguish three levels.

At the lowest level, a \textbf{digital model} exchanges data with its
physical counterpart only through manual steps: an engineer updates the
model when the plant changes, and nothing flows back automatically. This
is indistinguishable from a conventional CAD or CAE model and falls
short of twin status.

A \textbf{digital shadow} adds an automatic data flow in one direction,
from the physical object to the virtual one. The representation updates
without human intervention, but the virtual side cannot influence the
physical. This is still not a twin in the strict sense.

A \textbf{digital twin} completes the loop with bidirectional data flow:
the virtual model can send commands or schedules back to the physical
asset, and the twin is kept consistent with the physical system for the
length of its operating life. This bidirectional coupling is the
defining property that separates a twin from a shadow or a model.

These three levels are not merely terminological. They map onto a
deployment trajectory: a site typically begins with a digital shadow,
earns operator trust, and only then opens the feedback loop.

\hypertarget{synchronisation-strategies}{%
\subsection{Synchronisation
Strategies}\label{synchronisation-strategies}}

A digital twin is only as trustworthy as its last synchronisation. The
engineering question underneath every twin deployment is how to keep the
virtual state current without overwhelming the network or accepting
stale data. Chen and Devi {[}2{]} compare three synchronisation
strategies.

\textbf{Periodic pull} is the simplest: a central system polls the asset
at a fixed interval, regardless of what the asset is doing. The
advantage is implementation simplicity; the disadvantage is that quiet
assets are polled as often as busy ones, wasting bandwidth on unchanged
readings.

\textbf{Threshold push} reverses the initiative. The asset transmits
only when a sensed value moves beyond a configured delta. Bandwidth
follows activity, but a mis-set delta can hide slow drift entirely -- a
subtle failure mode because the twin appears live while its data is
stale.

\textbf{Event-driven flow} ties synchronisation to domain events: a job
start, a valve close, a product passing a sensor. The twin's state
history then reads as a process log rather than a time series. Operators
in Chen and Devi's trials strongly preferred this representation, and in
a synthetic benchmark across nine simulated production cells,
event-driven flows reduced transmitted volume by 71 per cent against
one-second periodic pull while keeping worst-case staleness bounded.

The choice of synchronisation strategy is not purely technical. It
shapes how operators read the twin's state and, by extension, whether
they trust it enough to act on it.

\hypertarget{factory-floor-applications}{%
\subsection{Factory-Floor
Applications}\label{factory-floor-applications}}

The factory floor is where digital twins move from pilots to production
value. Eriksen {[}3{]} documents an eighteen-month deployment of a
digital twin of a packaging line -- three filling stations, two
labellers, one palletiser -- at a mid-size beverage plant. The twin
began as a monitoring dashboard and ended as the plant's scheduling
authority.

The first six months deliberately stopped at a digital shadow: sensor
data flowed to the model automatically, and nothing flowed back. This
was an explicit trust-building decision. Operators could compare the
shadow against the line they could see, and its one visible product -- a
per-station utilisation display -- was checked against manual logs. Only
after the shadow was trusted did the deployment open the feedback loop.

The progression from shadow to twin is not unique to this case. It
reflects a broader pattern: sites that attempt bidirectional control
before earning trust tend to revert to monitoring, because operators
will not follow a twin's instructions if they do not believe the twin.

\hypertarget{interoperability-standards}{%
\subsection{Interoperability
Standards}\label{interoperability-standards}}

Digital-twin pilots rarely die of modelling problems; they die of
integration. Every asset on a real site speaks the protocol its vendor
shipped a decade ago, and the cost of a twin is dominated by the
adapters between those dialects and the twin's own model. Havel and
Ismail {[}4{]} argue that standards are how that cost is paid once
instead of per project.

Interoperability requires agreement at three distinct layers. At the
\textbf{conceptual layer}, the parties agree what the asset model
represents -- which physical components are in scope and what boundaries
the model draws. At the \textbf{semantic layer}, they agree what the
fields mean: that two systems' ``temperature'' is the same sensor, the
same unit, and the same sampling discipline. At the \textbf{transport
and syntax layer}, they agree how a message is structured -- a schema
language and a serialisation format.

Transport and syntax standards are mature and plentiful. Semantic
agreement is where projects stall, because it cannot be bought, only
negotiated. The practical instrument of semantic agreement is a shared
asset-model registry: a versioned catalogue mapping every signal to a
definition, a unit, and a provenance trail. When a new asset arrives on
site, its signals are looked up in the registry rather than
reverse-engineered from a vendor manual.

The three-layer model is not specific to digital twins, but it explains
a recurring pattern: twin projects that invest in semantic agreement
early survive integration, while those that treat it as a
post-deployment afterthought accumulate adapter debt that no amount of
modelling sophistication can offset.

\hypertarget{anomaly-detection-on-twin-state-streams}{%
\subsection{Anomaly Detection on Twin State
Streams}\label{anomaly-detection-on-twin-state-streams}}

Machine-learning anomaly detection falls outside this survey's scope,
but one source ties it directly to the twin and is worth noting briefly.
Farah and Gupta {[}5{]} observe that a digital twin's state stream is a
natural place to detect anomalies because it is already cleaned, aligned
and semantically labelled -- which constitutes most of the work in
industrial anomaly detection. The question their paper examines is how
minimal a detection model can be when it operates on a twin's state
stream rather than on raw sensor data. This is a narrow but instructive
example of how twin infrastructure can reduce the cost of downstream
analytics.

\hypertarget{comparison-of-approaches}{%
\subsection{Comparison of Approaches}\label{comparison-of-approaches}}

The following table summarises the papers surveyed, on the axis of what
each contributes to the twin literature.

\begin{longtable}[]{@{}
  >{\raggedright\arraybackslash}p{(\columnwidth - 6\tabcolsep) * \real{0.2500}}
  >{\raggedright\arraybackslash}p{(\columnwidth - 6\tabcolsep) * \real{0.2500}}
  >{\raggedright\arraybackslash}p{(\columnwidth - 6\tabcolsep) * \real{0.2500}}
  >{\raggedright\arraybackslash}p{(\columnwidth - 6\tabcolsep) * \real{0.2500}}@{}}
\caption{Where to start when building a first
twin.\label{tab:start-here}}\tabularnewline
\toprule\noalign{}
\begin{minipage}[b]{\linewidth}\raggedright
Starting point
\end{minipage} & \begin{minipage}[b]{\linewidth}\raggedright
Citekey
\end{minipage} & \begin{minipage}[b]{\linewidth}\raggedright
Core idea
\end{minipage} & \begin{minipage}[b]{\linewidth}\raggedright
Stated limitation
\end{minipage} \\
\midrule\noalign{}
\endfirsthead
\toprule\noalign{}
\begin{minipage}[b]{\linewidth}\raggedright
Starting point
\end{minipage} & \begin{minipage}[b]{\linewidth}\raggedright
Citekey
\end{minipage} & \begin{minipage}[b]{\linewidth}\raggedright
Core idea
\end{minipage} & \begin{minipage}[b]{\linewidth}\raggedright
Stated limitation
\end{minipage} \\
\midrule\noalign{}
\endhead
\bottomrule\noalign{}
\endlastfoot
Definition & {[}1{]} & Three-level taxonomy (model, shadow, twin) based
on integration depth & Taxonomy is illustrative; no empirical
validation \\
Synchronisation & {[}2{]} & Three synchronisation strategies compared;
event-driven preferred by operators & Synthetic benchmark across
simulated cells only \\
Factory deployment & {[}3{]} & Eighteen-month case study showing
shadow-to-twin progression & Single site; beverage packaging domain \\
Interoperability & {[}4{]} & Three-layer standards model; semantic
agreement as negotiation & Standards landscape simplified for
illustration \\
Anomaly detection & {[}5{]} & Twin state stream as natural
anomaly-detection input & Minimal model requirements unquantified \\
\end{longtable}

Table~\ref{tab:start-here}

The table reveals a pattern: every paper that focuses on a single
technical layer (synchronisation, interoperability, anomaly detection)
is strongest at that layer but does not address the others. The overview
paper {[}1{]} is the only one that attempts a cross-cutting taxonomy,
and its limitation -- illustrative rather than empirical -- is the
limitation of the field at large: few papers sit at the intersection of
definition, deployment, and standards.

\hypertarget{gaps}{%
\subsection{Gaps}\label{gaps}}

The retrieved corpus is small and covers the four scope themes, but
several gaps remain. First, there is no empirical comparison of
synchronisation strategies outside a synthetic benchmark; the
event-driven preference reported by Chen and Devi {[}2{]} has not been
validated on a live production line. Second, the interoperability
literature surveyed here is conceptual; there is no paper in the corpus
that measures the integration cost of a twin project before and after
adopting an asset-model registry. Third, the factory-floor case study by
Eriksen {[}3{]} is a single-site study in one domain; whether the
shadow-to-twin progression generalises to discrete manufacturing or
process industries is unknown.

These gaps are not failures of retrieval. The declared scope --
definitions, synchronisation, factory applications, and interoperability
-- is covered by the corpus. What the corpus does not contain is
empirical, multi-site validation of the patterns it describes, or a
quantitative study of integration cost as a function of standards
adoption. A later revision that widens the corpus should look for these
explicitly.

\hypertarget{references}{%
\subsection{References}\label{references}}

\hypertarget{refs}{}
\begin{CSLReferences}{0}{0}
\leavevmode\vadjust pre{\hypertarget{ref-sample_dt_overview_2024}{}}%
\CSLLeftMargin{{[}1{]} }%
\CSLRightInline{A. Author and B. Builder, {``Digital twins: Definitions,
distinctions, and a short taxonomy,''} \emph{Synthetic Sample Papers},
vol. 1, pp. 1--5, 2024.}

\leavevmode\vadjust pre{\hypertarget{ref-sample_dt_sync_2023}{}}%
\CSLLeftMargin{{[}2{]} }%
\CSLRightInline{C. Chen and D. Devi, {``State synchronisation strategies
for operational digital twins,''} \emph{Synthetic Sample Papers}, vol.
1, pp. 6--11, 2023.}

\leavevmode\vadjust pre{\hypertarget{ref-sample_dt_factory_2022}{}}%
\CSLLeftMargin{{[}3{]} }%
\CSLRightInline{E. Eriksen, {``A digital twin on the factory floor: An
eighteen-month case study,''} \emph{Synthetic Sample Papers}, vol. 1,
pp. 12--17, 2022.}

\leavevmode\vadjust pre{\hypertarget{ref-sample_std_interop_2021}{}}%
\CSLLeftMargin{{[}4{]} }%
\CSLRightInline{H. Havel and I. Ismail, {``Interoperability standards
for industrial asset models: A field guide,''} \emph{Synthetic Sample
Papers}, vol. 1, pp. 24--29, 2021.}

\leavevmode\vadjust pre{\hypertarget{ref-sample_ml_anomaly_2023}{}}%
\CSLLeftMargin{{[}5{]} }%
\CSLRightInline{F. Farah and G. Gupta, {``Lightweight anomaly detection
over digital-twin state streams,''} \emph{Synthetic Sample Papers}, vol.
1, pp. 18--23, 2023.}

\end{CSLReferences}

\end{document}
